% introduction_cn.tex - 引言

\section{引言}

面部微表情是个体试图隐藏或抑制真实情感时发生的非自愿、短暂面部运动\cite{ekman1969nonverbal}。与持续0.5--4秒的宏表情不同，微表情仅持续1/25至1/5秒\cite{ekman2003darwin}，具有低强度、部分面部参与和快速起始-顶点-偏移动态特性，即使对训练有素的编码者也难以检测和识别\cite{frank2009training}。

微表情识别的计算建模在过去十年取得显著进展，得益于CASME II\cite{yan2014casme2}、SAMM\cite{yap2018samm}、SMIC\cite{li2013smic}等基准数据集的发布。早期方法依赖手工时空特征，包括三正交平面局部二值模式（LBP-TOP）\cite{zhao2007lbptop}和光流描述符\cite{wang2015offapexnet}。然而这些方法受限于固定特征提取协议，无法从数据中学习层次化表示。

深度学习的出现催生了MER的范式转变。视频帧上的卷积神经网络、捕获时空体积的3D卷积网络\cite{tran2015c3d}以及利用自注意力机制的Transformer架构\cite{liu2022videoSwin}逐步推进了技术水平。尽管如此，现有深度MER系统存在关键局限：其架构对人类微表情感知的神经机制缺乏关注。

神经影像研究表明，面部表情感知——尤其是短暂或阈下表情——涉及大脑的双通路架构\cite{morris1999amygdala}。\textbf{快皮层下通路}从上丘经丘脑枕投射至杏仁核，实现对情感刺激的快速但粗糙处理。\textbf{慢皮层通路}涉及初级视皮层、梭状回面部区（FFA）和眶额皮层，支持精细辨别分析\cite{kanwisher1997ffa}。这两条通路并行运作并在杏仁核汇聚，调控注意力和情感反应。

本文提出\textbf{Censor}，一种明确模拟这一神经架构的仿生双通路框架。主要贡献如下：

\begin{enumerate}
\item \textbf{仿生架构设计}：提出双通路网络，包含处理光流的快通路3D ResNet-18（类比皮层下路径）和处理RGB+rPPG的慢通路3D Swin Transformer（类比皮层路径）。

\item \textbf{有效融合机制}：集成杏仁核启发注意力门控和挤压激励融合，结合两条通路的互补特征。

\item \textbf{混合专家分类}：噪声Top-2 MoE门控机制实现不同微表情类别的专家路由。

\item \textbf{全面实证验证}：在CASME II上使用留一受试者（LOSO）交叉验证进行评估，四类识别达到85\%准确率和0.80 F1分数。跨数据集实验证实了从CASME II迁移到SMIC和SAMM可获得0.74--0.76准确率的强泛化能力。
\end{enumerate}

本文结构如下：第2节回顾MER和仿生计算相关工作；第3节详细介绍Censor架构；第4节描述实验设置；第5节报告结果并进行讨论；第6节总结全文并展望未来工作。